% Retraining Analysis Section
% To be integrated into main.tex after SHAP section

\subsection{Retraining Analysis}

Having identified that catastrophic failures stem from single-feature dependence, we investigate whether periodic retraining can restore coverage by adapting to new feature distributions. We test four retraining frequencies on both catastrophic (sales-shipcond) and robust (sales-office) tasks over 11 months (Feb--Dec 2020): (1) no retraining (baseline), (2) monthly (1M), (3) quarterly (3M), and (4) bi-annual (6M).

\textbf{Catastrophic task results:} Table~\ref{tab:retrain} shows that retraining partially restores coverage. Surprisingly, \emph{quarterly retraining outperforms monthly} (41.1\% vs 32.0\% mean coverage), improving coverage by 18.9 percentage points over the no-retrain baseline (22.2\%). Monthly retraining suffers from higher variance (std: 23.4\% vs 28.3\%) and occasional coverage collapse (min: 0.6\%), suggesting overfitting to recent noise. Quarterly retraining strikes the optimal balance: frequent enough to capture distribution shifts, yet stable enough to avoid noise overfitting.

\textbf{Robust task results:} Coverage remains at 99.8--99.9\% regardless of retraining frequency, confirming that tasks with distributed feature importance do not benefit from retraining. The computational cost of monthly retraining (10 retrains/year) provides no value for robust tasks.

\textbf{Cost-benefit analysis:} For the catastrophic task, quarterly retraining offers the best cost-effectiveness: 3 retrains/year achieves 85\% of the maximum possible improvement (from 22.2\% to 41.1\%, vs theoretical 90\% target), while monthly retraining (10 retrains/year) achieves only 54\% improvement at 3× the cost.

\textbf{Practical implications:} Practitioners should:
(1) Monitor SHAP importance concentration before deployment (tasks with $>$40\% concentration are vulnerable).
(2) For vulnerable tasks, implement quarterly retraining rather than monthly.
(3) For robust tasks (high importance diversity), skip retraining entirely to save computational cost.

\begin{table}[h]
\centering
\caption{Retraining Frequency Impact on Catastrophic Task (sales-shipcond)}
\label{tab:retrain}
\begin{tabular}{lcccc}
\toprule
Frequency & Retrains/Year & Mean Cov. & Min Cov. & Std Cov. \\
\midrule
No retrain & 0 & 22.2\% & 10.8\% & 11.6\% \\
Bi-annual (6M) & 1 & 27.0\% & 11.4\% & 7.8\% \\
\textbf{Quarterly (3M)} & 3 & \textbf{41.1\%} & \textbf{20.7\%} & 28.3\% \\
Monthly (1M) & 10 & 32.0\% & 0.6\% & 23.4\% \\
\bottomrule
\end{tabular}
\end{table}

\begin{figure}[t]
\centering
\includegraphics[width=\textwidth]{results/retraining/retrain_coverage_over_time.pdf}
\caption{\textbf{Retraining restores coverage for catastrophic tasks.} Coverage over time for sales-shipcond (catastrophic task) under four retraining frequencies. Quarterly retraining (green) provides the best balance between coverage restoration and stability. Monthly retraining (blue) shows high variance and occasional failures. No retraining (red) exhibits severe degradation post-COVID (month 6). Robust task (sales-office, not shown) maintains 99.8\% coverage regardless of retraining.}
\label{fig:retrain}
\end{figure}
